% Item order of the solutions to Listas/pratica-lecture-05.md (25 items, in the list's order).
% \reuse{n}{slug} inputs a solution already written for lista-reais-lectures-01-05-solutions;
% \solitem{n}{slug} inputs a new sol-<slug>.tex of this folder.
\newcommand{\lecturevtitle}{Lecture 5 --- Repeated games}

\newcommand{\itemsv}{%
% Tier 1 --- Foundations
\reuse{1}{cabral-7-13}
\solitem{2}{cabral-7-8}
\solitem{3}{munoz-garcia-4-18}
\solitem{4}{gibbons-2-22}
\reuse{5}{tadelis-10-2}
\reuse{6}{munoz-garcia-6-1}
\solitem{7}{tadelis-10-9}
\solitem{8}{tadelis-9-3}
% Tier 2 --- Exam level
\solitem{9}{cabral-7-12}
\reuse{10}{tadelis-10-3}
\solitem{11}{gibbons-2-10}
\solitem{12}{tadelis-9-2}
\solitem{13}{tadelis-9-4}
\solitem{14}{munoz-garcia-6-6}
\reuse{15}{munoz-garcia-6-9}
\reuse{16}{gibbons-2-11}
\reuse{17}{tadelis-9-1}
\solitem{18}{tadelis-10-5}
\solitem{19}{tadelis-9-6}
% Tier 3 --- Advanced
\solitem{20}{mwg-12-aa-1}
\reuse{21}{tadelis-10-7}
\solitem{22}{tadelis-9-5}
\solitem{23}{gibbons-2-8}
\solitem{24}{gibbons-2-16}
\solitem{25}{mwg-12-aa-2}}

\newcommand{\solintro}{\sloppy These are detailed solutions to the 25 exercises of
\texttt{Listas/pratica-lecture-05.md} (Lecture~5, \emph{Repeated games}), in the list's order: three tiers (Foundations,
Exam level, Advanced) with the difficulty rising from 1 (apply a definition) to 5 (graduate level). Each heading shows the
item's number in the list, its level and its difficulty. Each statement is printed verbatim from the list right before its
solution. Each solution has the same parts: (a) what is asked; (b) the method, with the book section and page and the Lecture~5 slide;
(c) the solution, every step shown; (d) the economic interpretation; (e) the verification, by \texttt{check\_L05.py} in this folder
(new solutions) or \texttt{check\_C\_L05.py} in \texttt{lista-reais-lectures-01-05-solutions/} (the eight reused ones: items 1, 5, 6,
10, 15, 16, 17, 21), and against the book's worked answer where one exists; (f) common mistakes. Method of the slides: payoffs of
infinitely repeated games are normalised, $U_i = (1-\delta)\sum_t \delta^{t-1}u_i^t$; the one-stage deviation principle is checked history by history,
on and off the path; thresholds follow $\delta^* = (v^d - v^c)/(v^d - v^*)$ (Tadelis (10.5)). In the finitely repeated games (two stages) payoffs are
$u_i^1 + \delta u_i^2$. Three items deserve a warning: item~8 (Tadelis 9.3) has no solution for $\delta \le 1$ as printed, item~9
(Cabral 7.12) has no tension as printed, and item~20 (MWG 12.AA.1) is false as printed (the MWG manual's erratum); each solution says
what is true and why. Books: Tadelis (2013), \MG{} \& Toro-Gonzalez (2019, 2nd ed.), Gibbons (1992), Mas-Colell, Whinston \& Green (1995)
(MWG), and Cabral (2017). In bimatrices player~1 picks rows and the first payoff is player~1's; \underline{underlined} payoffs are
best replies. The eight reused solutions were written before the slides existed and cite ``Lecture~4, slides~21--22'' for the
repeated-game set-up; the matching slides of the real lecture are Lecture~5, slides~17--26.}
